% !TeX root = ../教材题库_全部条目_转写工作稿.tex
% Source: XJ-B1-C03-PR-P082-Y1193
% 第3章 / 3.2 函数的基本性质 / 教材 p.82 / PDF 第89页
% OCR未校对：公式与文字均可能有误；未完成数学TeX转写；禁止直接发给学生
% 原题及相邻原解截图（核对用）：../题目原图/XJ-B1-C03-PR-P082-Y1193-01.webp
% 本文件只在转写完成后作为题干使用；OCR 条目尚不是可靠题干。
\begin{bankitem}{p.82 练习 第2题}{第3章 / 3.2 函数的基本性质 / 教材 p.82 / PDF 第89页}{XJ-B1-C03-PR-P082-Y1193}{OCR未校对：公式与文字均可能有误；未完成数学TeX转写；禁止直接发给学生}
\begin{ocrdraft}
[教材p.82 / PDF第89页]
\par
2．检验下列函数的增减性，并说明是否有最大（小）值．如果有，指出最大
\par
（小）值和对应的最大（小）值点．
\par
2
\par
（2）不工）
\par
=3一
\par
2
\par
（3）不工）一一6十7()e[一2，4]）；
\par
（4）不工）一
\par
(xeCo,3]）．
\par
1十
\end{ocrdraft}
\end{bankitem}
